\documentclass{article}
\usepackage{graphicx} % Required for inserting images
\setlength{\parindent}{0pt}
\title{Lesson 5 Challenge - Calculus}
\usepackage{amsmath}
\author{William Wang}
\date{August 2026}

\begin{document}

\maketitle

\section{}
We want to prove that if we have 
\[\lim_{x \to 0} f(x) = L\]
then we have 
\[\lim_{x \to 0} f(cx)  = L\]
When we want to get the graph of \(f(cx)\) from \(f(x)\), we compress the graph horizontally by \(1/c\). So if the coordinates a point were \((a,b)\) then on the second graph it would be \((a/3, b)\). 

Then since we know that \(\lim_{x\to 0} f(x) = L\), by the definition of a limit, we have for every \(\epsilon > 0\) we have \(\delta_1 >0\) such that 
\[0 < |x| < \delta_1 \implies |f(x) - L|  < \epsilon\]
Let \(\epsilon > 0\), we need to find a \(\delta_2 > 0\) such if 
\[0 < |x| < \delta_2 \implies |f(cx) - L| < \epsilon\] Since \(c \neq 0\), we can define 
\[\delta_2 = \frac{\delta_1}{|c|}\]
\textbf{Proof}
By the definition of a limit, since 
\[\lim_{u\to0}f(u) = L\]
We know that an arbitrary \(\epsilon > 0\), there exists an \(\delta_1 > 0\) such that if
\[0 < |u| < \delta_1\]
then 
\[|f(u)-L| < \epsilon\]
Then let \(\epsilon > 0\) be given and we have \(\delta_2 = \frac{\delta_1}{|c|}\) and assume we have
\[0 < |x| < \delta_2 \]
then we can multiply the inequality by \(|c|\) to get 
\[0\cdot|c| < |c| \cdot|x| < \delta_2 \cdot|c|\]
\[0 < |cx| < \frac{\delta_1}{|c|}\cdot |c|\]
\[0 < |cx| < \delta_1\]
The let \(u = cx\). Since \(0 < |u| < \delta_1\), it follows that 
\[|f(u) - L| < \epsilon\]
Hence 
\[\lim_{x\to0} f(cx) =L\]
\section{}
Let 
\[f(x) = \begin{cases}
    3 & \text{if } x \geq 0 \\
    1  & \text{if } x < 0
\end{cases}\]
If we take the right hand limit, we have 
\[\lim_{x\to0^+} f(x) =3\]
since all the values we are using will be positive. However if we take the left hand limit
\[\lim_{x\to 0^-} f(x) = 1\]
since all values we are using will be negative. Since the definition of a limit is only fulfilled when the right hand limit and the left hand limit are the same, the limit doesn't exist. 

However if we have \(f(x^2)\) then all values will be positive since \(x^2 \ge 0\). Hence 
\[\lim_{x\to 0} f(x^2) = 3\]
Therefore the limit exists for \(f(x^2)\) but not \(f(x)\). 
\end{document}
